% page 466 du fac-simile (image p-465 du scan)
% bloc 154.9 x 242.0 mm ; marge gauche 8.0 mm
\pgc{-12.700mm}{0.000mm}{
\l{\cel{3}458}
\l{}
\l{\sou{prekara}. (Yorocienco).\cel{2}Grantita segun-requizite, e sempre}
\l{\cel{3}revokebla. - DEFIS.}
\l{}
\l{\sou{prekoca}. I. Di qua la matureso esas kresko-+anticipa. -}
\l{\cel{3}II. (\sou{metaf}.)Plu developita kam esas normala segun lua}
\l{\cel{3}evo. - EFIS.}
\l{}
\l{\sou{prekursoro}. Ta qua preparas la veno di altru. - EFIS.}
\l{}
\l{prelato. (\sou{eklezio katol.}) Oficiero ekleziala alta-grada. -}
\l{\cel{3}II. Oficiero ek la oficieraro papala, yurizita pri +werar}
\l{\cel{3}vesti violea. - DEFIRS.}
\l{}
\l{\sou{preliminaro}. Singla ek la elementi di to quo preiras e prepa-}
\l{\cel{3}ras la temo precipua. - DEFIRS.}
\l{}
\l{\sou{preludar}. (\sou{netrans.}) (\sou{muziko)}. Probar sua voco, sua muzik-}
\l{\cel{3}instrumento per kantar, plear serio de noti. - DEFIS.}
\l{}
\l{\sou{premio}. To quo donacesas kom rekompenso a persono qua vinkas}
\l{\cel{3}lor konkurso. - DEFIRS.}
\l{}
\l{\sou{premici}. La frukti dat-unesma di tero, di brutaro, quin onu}
\l{\cel{3}ofras a la deaji.\cel{2}- FIS.}
\l{}
\l{\sou{premiso}. (\sou{logiko}).Singla ek la du propozicioni (majoro e mi-}
\l{\cel{3}noro) di silogismo de qua onu inferas la konkluzo. (Uz-}
\l{\cel{3}esas precipue plurale). - DEFIS.}
\l{}
\l{\sou{prenar}. (\sou{trans., de,}\cel{1}ek). En-manuigar por tenar od uzar. FI}
\l{}
\l{\sou{preparar}. (\sou{trans., por}). Igar apta ye facar od agar ta a quo}
\l{\cel{3}lu esas destinita. - DEFIRS.}
\l{}
\l{\sou{preponderar}. (\sou{netrans., pri, en, che)}. Esar plu eminenta pri}
\l{\cel{3}autoritato, influo ; superesar (ulu, ulo) pri dominaco,}
\l{\cel{3}vinko. - DEFIS.}
\l{}
\l{\sou{prepoziciono}. (\sou{gram.}) Vorto-speco qua, lokizita avan nomo,}
\l{\cel{3}avan propoziciono infinitivala, e c., ligas ici per ra-}
\l{\cel{3}porto determinita a termino qua preiras. - DEFIRS.}
\l{}
\l{\sou{prepuco}. (\sou{anat.)}\cel{2}Prolonguro di la penis-pelo qua tegas la}
\l{\cel{3}glano. - EFIS.}
\l{}
\l{\sou{prerogativo}. Yuro atribuita ad ula rangi nobelesala privile-}
\l{\cel{3}jiera. - DEFIRS.}
\l{}
\l{\sou{presar}. (\sou{trans.}) Klemar per apogar su pezoze (sur lu). -}
\l{\cel{3}DEFIRS.}
\l{}
\l{presbita. (\sou{patol.)}\cel{1}(Persono) qua povas vidar nur to quo esas}
\l{\cel{3}for lu, pro ke lua muskulo ciliara cesis donar a la kris-}
\l{\cel{3}talino la kurviguro necesa. - DEFIS.}
}
